\newif\ifPDF
\ifx\pdfoutput\undefined\PDFfalse
\else\ifnum\pdfoutput > 0\PDFtrue
	\else\PDFfalse
	\fi
\fi

\ifPDF
	\documentclass[pdftex,11pt]{article}
	\RequirePackage[hyperindex,colorlinks,plainpages=false]{hyperref}
	\hypersetup{pdfauthor={Heng Li},linkcolor=blue,citecolor=blue,urlcolor=blue}
	\usepackage{graphicx}
	\DeclareGraphicsRule{*}{mps}{*}{}
\else
	\documentclass[11pt]{article}
	\usepackage{graphicx}
\fi

\newcommand{\pa}[2]{#1\otimes #2}

\addtolength{\textwidth}{3cm}
\addtolength{\hoffset}{-1.5cm}
\addtolength{\textheight}{4cm}
\addtolength{\voffset}{-2cm}

\makeindex

\title{Modelling Sequence Evolution by Stochastic Context-Free Grammar}
%\author{Heng Li and Richard Durbin}
\date{26 November 2006}

\usepackage{natbib}
\bibliographystyle{apalike}

\begin{document}

\maketitle

%\begin{abstract}
%  Stochastic context-free grammar (SCFG) is presented to model the
%  evolution of molecular sequences.  Fundamentally, the SCFG model is
%  almost identical to the classic likelihood model. However, working
%  with the inside-outside algorithm, SCFG smoothly and strictly gives
%  the posterior distributions of ancestral states and posterior
%  substitution matrix, which are usually complex to be calculated
%  directly in the classic likelihood framework. At the same time, SCFG
%  reveals the possibility to estimate parameters with expectation
%  maximization (EM). It also natrually implies the parsimony model when
%  the CYK algorithm is in use.
%\end{abstract}
%
%\section{Introduction}

%\subsection{ARG-CFG}
%I will not use strict language to describe the ARG-CFG. The model below
%can also be extended. I just show the model in a way that is best suited
%to ARG.
%
%Let $\mathcal{G}$ is an ARG constructed from $M$ people with $N$ sites.
%We label the recombination and coalescent nodes as $\{1,2,\ldots,n\}$.
%The root is always $n$.  Let $T^u$ be the marginal tree sliced from
%$\mathcal{G}$ at site $u$.  It should be noted that nodes should be
%present on the edges even if they have just one child.  For example,
%given an ARG $\tt (((1[m1],(2)5[r2])6,((5,3)7[m2],4)8[m3])9$, the
%marginal trees can be $\tt T^3=((1,(2)5)6,((3)7,4)8)9$ and $\tt
%T^1=T^2=((1)6,(((2)5,3)7,4)8)9$. This ARG corresponds to the example
%fig1 in Mark's paper, where $1=100$, $2=010$, $3=001$, $4=011$,
%$5=\mathrm{reombination}$. The first tree is fig1C, and the second is
%fig1B.  Sorry that I am still using my notation for ARG, because I think
%it is more concise.
%
%Define topological function $\delta^u_i(j,k)=1$ iff $j$ and $k$ are the
%two children of $i$ in tree $T^u$; otherwise $\delta^u_i(j,k)=0$. Also
%define $\delta^u_i(j)=1$ iff $j$ is the child of $i$ and $i$ is a
%single-child node (e.g. node 5 in $T^3$).
%$\delta^u_i(l)\cdot\delta^u_i(j,k)=0$ stands for all $i,j,k,l$ because
%$i$ can either be single-child or double-child.  Then we can define the
%SCFG.
%
%Nonterminal symbols: $[i,a,u]$, where $i$ is a node, $a\in\{0,1\}$ and
%$u$ is the site.  $S_0$ is the start non-terminal. Terminals are `(',
%`,', `)', `0' and `1'. Here are the production rules:
%\begin{center}
%\begin{tabular}{lll}
%  \hline
%  Type & Rule & Probability \\
%  \hline
%  Start $S_0$ & $S_0\rightarrow [n,a,u]$ & $q(a)r(u)$ \\
%  $i$ is single-child in $T^u$ & $[i,a,u]\rightarrow([j,b,u])$ & $\delta_i^u(j)p_{ij}(b|a)$ \\
%  $i$ is double-child in $T^u$ & $[i,a,u]\rightarrow([j,b,u],[k,c,u])$ & $\delta_i^u(j,k)p_{ij}(b|a)p_{ik}(c|a)$ \\
%  \hline
%\end{tabular}
%\end{center}
%This SCFG generates all the marginal trees and nothing else.

%In molecular phylogenetics, the three major inference methods are
%pairwise distances, maximum parsimony (MP) and the maximum likelihood
%(ML)~\citep{treereview-whelan01,book-felsenstein04}. Of the three
%methods, ML might be more frequently used in the literature, probably
%due to its great flexibility, versatility and robustness. ML inference
%not only permits the analyses of the evolutionary history with complex
%and precise models, but also opens a varity of possibilities for
%parameters estimation and statistical testing.  It has already achieved
%a great amount of meaningful results in phylogenetics.
%
%One of the intriguing topic in the ML framework is to map ancestral
%character states.  One groups of works focus on sophisticated sampling
%methods in a Bayesian
%framework~\citep{scfg-nielsen02,scfg-huelsenbeck03,scfg-pagel04,scfg-bollback06,scfg-ronquist04}.
%The other group solved this problem by theoretical analyses when they
%tried to implement the EM algorithms~\citep{em-dempster77} to estimate
%rate matrix and branch
%lengths~\citep{em-friedman02,em-holmes02,em-siepel04,scfg-klosterman06}.
%Possibly the second group thought the posterior ancestral states
%distribution was just an intermediate result. Their theory has not
%seemed to be widely used in mapping ancestral states.
%
%Theoretical deduction of the EM algorithm not only yields posterior
%ancestral states distribution, but also gives the posterior substitution
%matrices at each branch of a tree.  This was first implied by Friedman
%{\it at al.} in case of reversible evolution, and then futher developed
%independently by Holmes and Rubin, and by Siepel and Haussler.  Both are
%applicable to irreversible evolution. Interestingly, Holmes and Rubin
%symmetrized their substitution matrix at the last step. And thus their
%work was concluded to be best suited to reversible models by a more
%recent publication~\citep{scfg-klosterman06} which also established a
%theoretical framework for generalized rate matrix, including
%substitution matrix, estimations.
%
%In this paper, we recovered the previous results in calculating
%posterior distribution of ancestral states and substitution matrices at
%each branch.  These results are achieved with a stochastic context-free
%grammar (SCFG)~\citep{durbin98}, which resembles the current likelihood
%model in many ways, but brings new tools in favour of the inside-outside
%algorithm, an algorithm native to the SCFG theory.
%
\section{Methods and Results}

\subsection{Brief introduction to stochastic context-free grammar}
Context-free grammar (CFG) provides a framework for modelling some
strings of symbols.  First developed by computational
linguists~\citep{scfg-chomsky56}, CFGs, with long range dependencies such
as parentheses, have come to play an important role in computer science,
for example for designing compilers for computer languages. In the area
of bioinformatics, CFGs are mainly used to model the secondary
structures of RNAs~\citep{scfg-eddy94,durbin98}. In this paper, we will
find its another application in modelling sequence evolution.

A CFG is a type of transformational grammar. A transformation grammar is
comprised of four components: a set of terminal symbols which can be
observed, a set of nonterminal symbols which are hidden and can be
derived, a set of production rules which define the transformations, and
a start nonterminal symbol from which all strings are generated.  In a
transformational grammar, production rules are presented in a form like
$\alpha\rightarrow\gamma$, where $\alpha$ is a string containing at
least one nonterminal symbol, and $\gamma$ is any string.  The
definition of a CFG further requires that there should be only one
symbol in left-hand side of each rule.  A rule like $SS\rightarrow S$ is
allowed in an unrestricted transformational grammar, but not in a CFG.

An example of a CFG that can generate a palindrome language is:
\[
S\rightarrow aSa|bSb|aa|bb
\]
In this example, $\{S\}$ is the nonterminal set, $\{a,b\}$ is the
terminal set, $S$ is the start symbol, and $\{S\rightarrow
aSa,S\rightarrow bSb,S\rightarrow aa,S\rightarrow bb\}$ is the set of
production rules.  A string like $baabbaab$ can be generated in the
following procedure:
\[
S\Rightarrow bSb\Rightarrow baSab\Rightarrow baaSaab\Rightarrow baabbaab
\]
This procedure is called a derivation and will be denoted by $\pi$ in
this paper.  A derivation is a series of consecutive rules deriving from
the start state to a string of terminal states.

In the above example, $baabbaab$ correponds to one derivation. However,
if we add a new rule `$S\rightarrow SS$', the following derivation can
also result in the same string:
\[
S\Rightarrow SS\Rightarrow bSbbSb\Rightarrow baabbaab
\]
Which derivation is `preferred'? In fact, we can assign a probability to
each production rule and require each nonterminal symbol is transformed
with a probability $1.0$.  Assuming rules are applied independently
(Markov property), we are able to find the most probable derivation with
the CYK algorithm, to calculate the probabilities of strings or
derivations with the inside or outside algorithm, and to estimate the
parameters or their posterior distribution with the inside-outside
algorithm.  A CFG with probabilities assigned in this way is called a
stochastic context-free grammar, or an SCFG in brief.

\subsection{The construction of the SCFG}
Let $V$ be a set of leaves, $\mathcal{T}$ be a set of {\it binary
  rooted} trees whose leaf set is $V$. We consider two internal nodes
are the same if and only if the sets of leaves descend from the two
nodes are identical. Then we can label all the nodes as
$\mathcal{U}(\mathcal{T})=\{1,2,\ldots,n\}$. For convenience, the root
node, which will appear in each tree in $\mathcal{T}$, is always labeled
as $n$.

To describe the topologies in $\mathcal{T}$, we introduce a function
$\delta_i(j,k)$, $1\leq i,j,k\leq n$.  The value of this function is not
zero only if: i) node $j$ and $k$ are $i$'s children in a certain tree
$T\in\mathcal{T}$ and ii) $j\leq k$. We can give $\delta_i(j,k)$ a
probabilistic interpretion by requring $\sum_{j,k}\delta_i(j,k)=1$.

Let $\Sigma$ be an alphabet. Its elements are called
characters. Practically, $\Sigma$ may be the set of the four
nucleotides, of the 20 amino acids, or of present/absent states of a
marker.

Based on $\mathcal{T}$ and $\Sigma$, an SCFG can be constructed to
generate a binary rooted tree together with the characters assigned to
the leaves.  Except the start symbol $S_0$, a nonterminal symbol of the
SCFG is a pair $[i,a]\in \mathcal{U}(\mathcal{T})\times\Sigma$; a
terminal symbol is $a\in\Sigma\cup\{\mbox{`(',`)',`,'}\}$. The
production rules and their probabilities are presented below:
\begin{center}
\begin{tabular}{lll}
  \hline
  Type & Rule & Probability \\
  \hline
  start symbol & $S_0\rightarrow[n,a]$ & $q(a)$ \\
  $i$ is internal & $[i,a]\rightarrow([j,b],[k,c])$ & $p_{ij}(b|a)p_{ik}(c|a)\delta_i(j,k)$ \\
  $i$ is external & $[i,a]\rightarrow v$ & $1$ \\
  \hline
\end{tabular}
\end{center}
where $p_{ij}(b|a)$ is the probability that character $a$ is transformed
to $b$ on branch $(i,j)$, and $q(a)$ is the probability of character $a$
appearing at the root. If we require \mbox{$\sum_a q(a)=1$}, and $\sum_b
p_{ij}(b|a)=1$, all the nonterminal symbols will be derived with a
probability $1.0$. The grammar given above is an SCFG, an SCFG that
generates a string describing a tree over leaves in the New Hampshire
format.

To help the description of algorithms, we write $T_i$, $1\leq i\leq n$,
as the subtree descended from node $i$, and write $\sigma(T_i)$ as the
New Hampshire string for $T_i$.  A notation like $\sigma_1$ or
$\sigma_2$ always denotes a string.

\subsection{The inside algorithm}
In the SCFG framework, the inside algorithm is used to calculate the
probability of a string being generated from a symbol. In our particular
case, it is $\Pr\{[i,a]\Rightarrow\sigma(T_i)\}$, the probability of a
tree string being generated from $[i,a]$. This probability $\alpha(i,a)$
can be recursively calculated as follows:
\begin{eqnarray}\label{equ:in}
  \alpha(i,a)&=&\Pr\{[i,a]\Rightarrow\sigma(T_i)\}\nonumber\\
  &=&\sum_{j,k}\sum_{b,c}\Pr\{[i,a]\Rightarrow([j,b],[k,c])\Rightarrow\sigma(T_j)\sigma(T_k)\}\nonumber\\
  &=&\sum_{j,k}{\delta_i(j,k)\sum_b{\alpha(j,b)p_{ij}(b|a)}\sum_c{\alpha(k,c)p_{ik}(c|a)}}
\end{eqnarray}
Then the probability of generating $\mathcal{T}$ from $S_0$ is:
\begin{eqnarray}
  P(\mathcal{T})&=&\Pr\{S_0\Rightarrow\sigma(\mathcal{T})\}\nonumber\\
  &=&\sum_v q(a)\Pr\{S_0\Rightarrow[n,a]\Rightarrow\sigma(\mathcal{T})\}\nonumber\\
  &=&\sum_v q(a)\alpha(n,a)
\end{eqnarray}
The inside algorithm, as well as the CYK algorithms, works in a
bottom-up way.  In calculation, children are always calculated earlier
than their parents; $\alpha(i,a)$ for leaves will be set first.

In case where $\mathcal{T}=\{T\}$, the inside algorithm for this SCFG is
identical to Felsenstein's {\it pruning}
algorithm~\citep{felsenstein81a}.

\subsection{The posterior expectation of ancestral character states}

The inside algorithm calculates the probability of a tree. However, in
order to calculate the posterior expectation, we also need to calculate
some probabilities in a top-down way. This is the outside algorithm.
The equivalent algorithm has already been described in previous
works~\citep{em-holmes02,em-siepel04}.

The outside algorithm calculates the probability of a part of derivation
which stops at a symbol $[i,a]$. For the root node $n$, the calculation
is simple:
\begin{eqnarray}
\beta(n,a)=\Pr\{S_0\Rightarrow[n,a]\}=q(a)
\end{eqnarray}
For $i<n$, this probability $\beta(i,a)$ can also be recursively calculated:
\begin{eqnarray}\label{equ:out}
  \beta(i,a)&=&\Pr\{S_0\Rightarrow\sigma_1[i,a]\sigma_2\}\nonumber\\
  &=&\Pr\{S_0\Rightarrow\sigma_1[i,a]\sigma_2,\mbox{$i$ is the left child}\}\nonumber\\
  &&+\Pr\{S_0\Rightarrow\sigma_1[i,a]\sigma_2,\mbox{$i$ is the right child}\}\nonumber\\
  &=&\sum_{j,k}\sum_{b,c}\Pr\{S_0\Rightarrow\sigma_1[k,c]\sigma'_2\Rightarrow\sigma_1([i,a],[j,b])\sigma'_2\Rightarrow\sigma_1[i,a]\sigma(T_j)\sigma'_2\}\nonumber\\
  &&+\sum_{j,k}\sum_{b,c}\Pr\{S_0\Rightarrow\sigma'_1[k,c]\sigma_2\Rightarrow\sigma'_1([j,b],[i,a])\sigma_2\Rightarrow\sigma'_1\sigma(T_j)[i,a]\sigma_2\}\nonumber\\
  &=&\sum_{j,k}\delta_k(j,i)\sum_{b,c}{\alpha(j,b)\beta(k,c)p_{kj}(y|z)p_{ki}(v|z)} \nonumber \\
  &&+\sum_{j,k}\delta_k(i,j)\sum_{b,c}{\alpha(j,b)\beta(k,c)p_{ki}(v|z)p_{kj}(y|z)}
\end{eqnarray}
Note that given a single tree, node $i$ can either be the left or the
right child, but cannot be the both.  As a result, one of
$\delta_k(j,i)$ and $\delta_k(i,j)$ must equal to zero.

The combination of inside and outside algorithms, usually termed as the
inside-outside algorithm, provides a way to calculate the posterior
expectation and variance of some statistics.

Let $X_{ijk,abc}(\pi)$ be the number of times when the rule
$[i,a]\rightarrow([j,b],[k,c])$ is in use given a derivation $\pi$. A
derivation $\pi$ is a random variable, and so it is with
$X_{ijk,abc}(\pi)$.  In a single derivation of this SCFG, each rule will
be used only once or not be used at all, so $X_{ijk,abc}$ will be either
0 or 1. Thus the calculation of posterior expectation $x_{ijk,abc}$ of
$X_{ijk,abc}$ is reduced to the calculation of
$\Pr\{X_{ijk,abc}=1|S_0\Rightarrow\sigma(T)\}$:
\begin{eqnarray}\label{equ:c1}
  x_{ijk,abc}&=&\sum_{\pi}X_{ijk,abc}(\pi)P(\pi|\sigma(T))\nonumber\\
  &=&\Pr\{S_0\Rightarrow\sigma_1[i,a]\sigma_2\Rightarrow\sigma_1([j,b],[k,c])\sigma_2\Rightarrow\sigma(T)|S_0\Rightarrow\sigma(T)\} \nonumber\\
  &=&\frac{1}{P(\sigma(T))}\delta_i(j,k)\beta(i,a)\alpha(j,b)p_{ij}(b|a)\alpha(k,c)p_{ik}(c|a)
\end{eqnarray}
As $X_{ijk,abc}$ can be either 0 or 1, $X^2_{ijk,abc}=X_{ijk,abc}$
always stands. This makes it possible to calculate the variance of
$X_{ijk,abc}$:
\begin{equation}
v_{ijk,abc}={\rm Var}X_{ijk,abc}=x_{ijk,abc}[1-x_{ijk,abc}]
\end{equation}

So far, we have only considered the situation that one leaf has just one
site.  However, in case of a tree built from a multi-sequence alignment,
there are usually many of sites. We denote the character string
corresponding to site $u$ by $\sigma^{(u)}(T)$, and the sub-derivation
corresponding to site $u$ by $\pi^{(u)}$. We write $\vec{\sigma}(T)$ as
the character string of all sites. Similarly we have $\vec{\pi}$, the
entire derivation for all sites. If we assume each site is independent
of all the others, the expectation of
$\tilde{X}_{ijk,abc}(\vec{\pi})=\sum_{\vec{\pi}}X^{(u)}_{ijk,abc}(\pi^{(u)})$
is
\begin{eqnarray}\label{equ:cijk-all}
  \tilde{x}_{ijk,abc}&=&\sum_{\vec{\pi}}\tilde{X}_{ijk,abc}(\vec{\pi})P(\vec{\pi}|\vec{\sigma}(T))\nonumber\\
  &=&\sum_{\vec{\pi}}\tilde{X}_{ijk,abc}(\vec{\pi})\prod_u P(\pi^{(u)}|\sigma^{(u)}(T))\nonumber\\
  &=&\sum_u\sum_{\pi^{(u)}}X^{(u)}_{ijk,abc}(\pi^{(u)})P(\pi^{(u)}|\sigma^{(u)}(T))\nonumber\\
  &=&\sum_u x^{(u)}_{ijk,abc}
\end{eqnarray}
where $x^{(u)}_{ijk,abc}$ is calculated by Equation~\ref{equ:c1}. In
addition, using the site independence again, we will get the variance in
case of multiple sites:
\begin{equation}\label{equ:cijk-var-all}
  \tilde{v}_{ijk,abc}={\rm Var}\tilde{X}_{ijk,abc} = \sum_u x^{(u)}_{ijk,abc}\Big[1-x^{(u)}_{ijk,abc}\Big]
\end{equation}
Note that when the number of sites is sufficiently large, the posterior
distribution of $\tilde{X}_{ijk,abc}$ can be approximated by a normal
distribution $N(\tilde{x}_{ijk,abc}, \tilde{v}_{ijk,abc})$.

We next introduce the marginal expectation $\tilde{x}_{ij,ab}$,
$\tilde{x}_{i,a}$ and $\tilde{x}_{ijk}$ for further use:
\begin{equation}
  \tilde{x}_{ij,ab}=\left\{\begin{array}{ll}
      \sum_{k,c}\tilde{x}_{ijk,abc} & \mbox{if $j$ is the left child}\\
      \sum_{k,c}\tilde{x}_{ikj,azy} & \mbox{if $j$ is the right child}\\
\end{array}\right.
\end{equation}
\begin{equation}
  \tilde{x}_{i,a}=\sum_{j,b}\tilde{x}_{ij,ab}
\end{equation}
\begin{equation}
  \tilde{x}_{ijk}=\sum_{a,b,c}\tilde{x}_{ijk,abc}
\end{equation}
\begin{equation}
  \tilde x_{i}=\sum_{j,k}\tilde{x}_{ijk}=\sum_v\tilde{x}_{i,a}
\end{equation}
Note that $\tilde{x}_{i,a}$ is the posterior expectation of
$\tilde{X}_{i,a}$, the number of of times when a rule $[i,a]$ is
used. Applying Equation~\ref{equ:in}, we have:
\begin{equation}
\tilde{x}_{i,a}=\sum_{j,k,b,c}x_{ijk,abc}=\sum_u\frac{\alpha(i,a)\beta(i,a)}{P(\sigma(T))}
\end{equation}


\subsection{Parameter estimation by expectation maximization}
The calculation of the posterior expectation of $\tilde{X}_{ijk,abc}$
paves the way for the application of the EM
algorithm~\citep{em-dempster77} to estimate the SCFG parameters.  Based
on a known set of parameters $\theta^t$, the EM algorithm estimate a new
set of parameters $\theta$ which will increase
$P(\sigma(T)|\theta)$. This process is usually iterated until a local
maximum is reached (in some case this can be shown to be global).
Throughout this section, notations like $\tilde{x}^t_i(v)$ indicates
that $\tilde{x}^t_{i,a}$ is calculated with $\theta^t$, while
$\tilde{x}_{i,a}$ with $\theta$.

Given a derivation $\vec{\pi}$, we know that
\begin{equation}
  P(\vec{\sigma}(T),\vec{\pi}|\theta)=\prod_{i,\delta_i(j,k)=1}\prod_{a,b,c}\Big[p_{ij}(b|a)p_{ik}(c|a)
  \delta_i(j,k)\Big]^{\tilde{X}_{ijk,abc}(\vec{\pi})}
  \prod_u q(d)^{\tilde{X}_{n,d}(\vec{\pi})}
\end{equation}
Then the $Q$-function in the EM iteration is:
\begin{eqnarray}\label{equ:em-q}
  Q(\theta|\theta^t)&=&\sum_{\vec{\pi}}P(\vec{\pi}|\vec{\sigma}(T),\theta^t)\log P(\vec{\sigma}(T),\vec{\pi}|\theta)\nonumber\\
  &=&\sum_{\vec{\pi}}P(\vec{\pi}|\vec{\sigma}(T),\theta^t)\Bigg\{\sum_{i,j,k}\sum_{a,b,c}\tilde{X}_{ijk,abc}(\vec{\pi})\Big[\log p_{ij}(b|a)+\log p_{ik}(c|a)\nonumber\\
  &&+\log\delta_i(j,k)\Big]+\sum_d \tilde{X}_{n,d}(\vec{\pi})\log q(d)\nonumber\Bigg\}\\
  &=&\sum_{i,j}\sum_{a,b}\tilde{x}^t_{ij,ab}\log p_{ij}(b|a)+\sum_{i,j,k}\tilde{x}_{ijk}\log\delta_i(j,k)
  +\sum_d \tilde{x}^t_{n,d}\log q(d)
\end{eqnarray}
Based on the theory of relative entropy~\citep{durbin98}, we know that
$Q(\theta|\theta^t)$ is maximized when $q(d)=\hat{q}(d)$,
$p_{ij}(b|a)=\hat{p}_{ij}(b|a)$ and $\delta_i(j,k)=\hat{\delta}_i(j,k)$,
where:
\begin{eqnarray}
\hat{p}_{ij}(b|a)&=&\frac{\tilde{x}^t_{ij,ab}}{\tilde{x}^t_{i,a}}\\
\hat{\delta}_i(j,k)&=&\frac{\tilde{x}^t_{ijk}}{\tilde{x}^t_i}\\
\hat{q}(d)&=&\frac{\tilde{x}^t_{n,d}}{\tilde{x}^t_n}
\end{eqnarray}
Iterating the calculation of $q(d)$, $\delta_i(j,k)$ and $p_{ik}$ will
give the EM estimation of all the parameters in the SCFG model.  We may
also use a symmetric substitution matrix, or use a unique matrix across
all the branches. The iteration equations in those cases can also be
analytically deduced in a similar way.

\subsection{The CYK algorithm}
For the sake of theoretical completeness, we present the CYK algorithm
as the last one.  The CYK algorithm finds the most probable derivation
given a string of terminal symbols.  Define $\gamma(i,a)$ as:
\begin{equation}
  \gamma(i,a)=\Pr\{\mbox{most probable procedure to generate $\sigma(T_i)$ from $(i,a)$}\}
\end{equation}
Then $\gamma(i,a)$ can be calculated by a straightforward iteration:
\begin{eqnarray}
  \gamma(i,a)&=&\max_{j,k,b,c}\Big\{\gamma(j,b)\gamma(k,c)\Pr\{[i,a]\rightarrow([j,b],[k,c])\}\Big\}\nonumber\\
  &=&\max_{j,k,b,c}\Big\{\delta_i(j,k)\gamma(j,b)\gamma(k,c)p_{ij}(b|a)p_{ik}(c|a)\delta_i(j,k)\Big\}
\end{eqnarray}

It is worth noting that the CYK algorithm is reduced to the weighted
parsimony if we use a unique substitution matrix across all the
branches. The weight of a pair of characters $(a,b)$ is `$-\log
p(b|a)$'. This finding agrees with the probabilistic interpretion of
parsimony~\citep{felsenstein81b}.

Given a tree, the CYK algorithm works in a bottom-up way. However, in
our case, we can also apply it in a top-down way without any tree. This
will give us the most probable tree among all trees that can be
generated from the SCFG. We call the top-down CYK as the tree merge
algorithm.

\subsection{The tree merge algorithm}

\section{Discussions}

\subsection{Time complexity}
The time complexity of constructing $\mathcal{U}(\mathcal{T})$ is
$O(|\mathcal{T}|\cdot|V|^2)$.  This includes traversing all the nodes
and set comparisons with hash. Parsing one tree with inside or outside
algorithm can be done in $O(|\mathcal{U}|\cdot|\Sigma|^2)$. This is true
because for one tree $\sum_{j,k}$ in Equation~\ref{equ:in} and
\ref{equ:out} are not carried out. In the next step, directly
calculating all $x_{ijk,abc}$ needs $O(|\mathcal{U}|\cdot|\Sigma|^3)$.
However, if we notice that only $x_{ij,ab}$ and $x_{ijk}$ are
practically useful, we can achieve this step in
$O(|\mathcal{U}|\cdot|\Sigma|^2)$, saving a $|\Sigma|$.  Usually we are
dealing with $L$ sites. The inside-outside algorithm should then be
repeated for $L$ times.  As a consequence, up to the calculation of
posterior expectation, the time complexity is
$O(|\mathcal{T}|\cdot|V|^2+L|\mathcal{U}|\cdot|\Sigma|^2)$ in
total. When we use EM, the number of iterations should be multiplied to
the second term.


In this paper, we present a stochastic context-free grammar that is
compatible with both maximum likelihood and maximum parsimony
methods. As an extension to the two classic methods, this SCFG also
gives the posterior distribution of ancestral state with
Equation~\ref{equ:civ}, and the posterior expectation and variance of
number of substitutions with Equation~\ref{equ:cijk-all} and
\ref{equ:cijk-var-all}. It shows that these results can be calculated
directly without complex simulation processes.

This SCFG also reveals the possibility for parameter estimation with the
EM algorithm.  Although for a common parameter estimation problem we can
hardly get analytical solution from Equation~\ref{equ:em-q}, this
equation does show another promissing aspect: branch-specific parameters
may be separated in optimization. As a matter of fact, the first sum in
Equation~\ref{equ:em-q} is calculated over all branches. If each
$p_{ij}$ is only determined by branch-specific parameters, such as
branch lengths, the global optimization of $Q(\theta|\theta^t)$ can be
reduced to a series of sub-optimizations at each branch. This will
effectively avoid high-dimension multi-variable optimization which
usually consumes most of CPU time in ML estimation. Classic ML methods
attempt to reduce the computaional burden by using less linage specific
parameters, while SCFG-EM algorithm may be more efficient if all the
parameters are linage specific.

Context-free grammar has a natrual connection with a tree topology. In
this paper, we only show an SCFG that generates the string of symbols of
leaves, but in general, SCFG can be used to generate a string in Newick
format, and thus describes the full topology of a tree.  We also know
that SCFG can model macro-gene evolution and mimic the majority rule
consensus tree algorithm. Detailed discussion of these applications may
be beyond the scope of this paper.

\bibliography{ref}

\end{document}
